\documentclass[letterpaper]{report}

\usepackage[T1]{fontenc}
\usepackage{helvet}
\renewcommand{\familydefault}{\sfdefault}

\usepackage[bookmarks=true,bookmarksnumbered=true]{hyperref}
\usepackage{bookmark}

% Make generated PDFs reproducible across machines and build times.
\ifdefined\pdfinfoomitdate
\pdfinfoomitdate=1
\fi
\ifdefined\pdftrailerid
\pdftrailerid{}
\fi
\ifdefined\pdfsuppressptexinfo
\pdfsuppressptexinfo15
\fi
% Format chapter headings as "1 Introduction" (no "Chapter" prefix)
\usepackage{titlesec}
\titleformat{\chapter}[hang]{\huge\bfseries}{\thechapter}{1em}{}

% Remove page numbers to match the original
\pagestyle{empty}
\usepackage{etoolbox}
\patchcmd{\chapter}{\thispagestyle{plain}}{\thispagestyle{empty}}{}{}

\begin{document}

\chapter{Collective Communication}

The simplest collective communication function is \texttt{MPI\_Barrier}, which blocks
until all processes in the communicator have reached the barrier.

\begin{verbatim}
int MPI_Barrier(MPI_Comm comm)
\end{verbatim}

The broadcast function \texttt{MPI\_Bcast} sends data from one process to all others.

\begin{verbatim}
int MPI_Bcast(void* buffer, int count, MPI_Datatype datatype,
int root, MPI_Comm comm)
\end{verbatim}

\chapter{More Collective Operations}

The gather function \texttt{MPI\_Gather} collects data from all processes.

\begin{verbatim}
int MPI_Gather(const void* sendbuf, int sendcount, MPI_Datatype sendtype,
void* recvbuf, int recvcount, MPI_Datatype recvtype,
int root, MPI_Comm comm)
\end{verbatim}

The scatter function distributes data from one process to all others.

\begin{verbatim}
int MPI_Scatter(const void* sendbuf, int sendcount, MPI_Datatype sendtype,
void* recvbuf, int recvcount, MPI_Datatype recvtype,
int root, MPI_Comm comm)
\end{verbatim}

\end{document}
